\subsection{Арифметические операторы}

\textbf{Мотивирующий вопрос:} Как научить дрона считать расстояние или время полета?

Python поддерживает все стандартные математические операции.

\begin{itemize}
    \item \texttt{+} Сложение
    \item \texttt{-} Вычитание
    \item \texttt{*} Умножение
    \item \texttt{/} Деление (всегда возвращает \texttt{float})
    \item \texttt{//} Целочисленное деление (отбрасывает дробную часть)
    \item \texttt{\%} Остаток от деления
    \item \texttt{**} Возведение в степень
\end{itemize}

\begin{codefile}[python]{examples/python/03_arithmetic.py}
\end{codefile}

Обратите внимание на приоритет операций: умножение и деление выполняются раньше сложения и вычитания, как в математике. Скобки \texttt{()} используются для изменения порядка действий.

\begin{taskbox}[Калькулятор]
Напишите программу, которая считает площадь прямоугольника со сторонами 5 и 8 метров.
\end{taskbox}
